\documentclass[hyperref={pdfpagelabels=false},table,handout]{beamer}

\input{lec_common}
\begin{document}

\part{Unit 3: More on Lists and other Container Types}


\begin{frame}[fragile]{Creating sublists}
\begin{block}{Slicing -- \pyth!L[i:j]!}
\alert{Slicing} a list between index $i$ and $j$ is forming a new list by taking
elements with indices starting at $i$ and ending \emph{just before} $j$.
\end{block}

\begin{block}{Alternatively ...}
\pyth!L[i:j]! means: create a list by removing the first
$i$ elements from $L$  and keeping the next $j-i$ elements (for $j > i \ge 0 $)
and removing the rest.
\end{block}
\vfill
\begin{example}
\begin{python}
L = ['C', 'l', 'o', 'u', 'd']
L[1:4] # remove one element and take three from there:
# ['l', 'o', 'u']
\end{python}
\end{example}
\end{frame}

\begin{frame}[fragile]{Partial slicing}
One may omit the first or last bound of the slicing:
\begin{python}
L = ['C', 'l', 'o', 'u', 'd']
L[1:] # ['l', 'o', 'u', 'd']
L[:3] # ['C', 'l', 'o']
L[-2:] # ['u', 'd']    # negative indexing !
L[:-2] # ['C', 'l', 'o']
L[:] # the entire list
\end{python}
\vfill
\begin{block}{Mathematical Analogy}
This is similar to half lines in $\mathbb{R}$:
$[-\infty,a)$ means: take all numbers strictly lower than $a$,
cf.  syntax of  \pyth!L[:j]!.
\end{block}
Let's sum it up
\begin{itemize}
\item \pyth!L[i:]! take all elements \emph{except} the $i$ first ones
\item \pyth!L[:i]!  take the $i$ first elements
\item \pyth!L[-i:]!  take the last $i$ elements
\item \pyth!L[:-i]! to take all elements \emph{except} the $i$ last ones
\end{itemize}
\end{frame}

%\tikzstyle{filled}=[fill=green!10]
%\tikzstyle{empty}=[fill=red!80]
\tikzstyle{filled}=[fill=red!80]
\tikzstyle{empty}=[fill=green!10]

\begin{frame}[fragile]{Examples}
\begin{itemize}
\item \pyth!L[2:]! \hfill
\begin{tikzpicture}
\matrix[nodes={draw,style=filled},minimum size=0.5cm,ampersand replacement=\&]
{
\node[style=empty]{0} ; \& \node[style=empty]{1}; \& \node{2};  \& \node{3}; \& \node{$\cdots$}; \& \node{-3}; \& \node{-2}; \& \node{-1}; \\
};\end{tikzpicture}
\item \pyth!L[:2]! \hfill
\begin{tikzpicture}
\matrix[nodes={draw,style=empty},minimum size= 0.5cm,ampersand replacement=\&]
{
\node[style=filled]{0} ; \& \node[style=filled]{1}; \& \node{2};  \& \node{3}; \& \node{$\cdots$}; \& \node{-3}; \& \node{-2}; \& \node{-1}; \\
};
\end{tikzpicture}
\item \pyth!L[:-2]! \hfill
\begin{tikzpicture}
\matrix[nodes={draw,style=filled},minimum size= 0.5cm,ampersand replacement=\&]
{
\node{0} ; \& \node{1}; \& \node{2};  \& \node{3}; \& \node{$\cdots$}; \& \node{-3}; \& \node[style=empty]{-2}; \& \node[style=empty]{-1}; \\
};
\end{tikzpicture}
\item \pyth!L[-2:]! \hfill
\begin{tikzpicture}
\matrix[nodes={draw,style=empty},minimum size= 0.5cm,ampersand replacement=\&]
{
\node{0} ; \& \node{1}; \& \node{2};  \& \node{3}; \& \node{$\cdots$}; \& \node{-3}; \& \node[style=filled]{-2}; \& \node[style=filled]{-1}; \\
};
\end{tikzpicture}
\item \pyth!L[2:-1]! \hfill
\begin{tikzpicture}
\matrix[nodes={draw,style=filled},minimum size= 0.5cm,ampersand replacement=\&]
{
\node[style=empty]{0}; \& \node [style=empty]{1}; \& \node{2};  \& \node{3}; \& \node{$\cdots$}; \& \node{-3}; \& \node{-2}; \& \node [style=empty]{-1}; \\
};
\end{tikzpicture}
\item \pyth!L[2:5]! \hfill
\begin{tikzpicture}
\matrix[nodes={draw,style=empty},minimum size= 0.5cm,ampersand replacement=\&]
{
\node{0}; \& \node{1}; \& \node[style=filled]{2};  \& \node[style=filled]{3}; \& \node[style=filled]{4}; \& \node{5}; \& \node{$\cdots$}; \& \node{-1}; \\
};
\end{tikzpicture}
\item \pyth!L[-4:-1]! \hfill
\begin{tikzpicture}
\matrix[nodes={draw,style=empty},minimum size= 0.5cm,ampersand replacement=\&]
{
\node{0}; \& \node{1}; \& \node{$\cdots$};  \& \node{-5}; \& \node[style=filled]{-4}; \& \node [style=filled]{-3}; \& \node [style=filled]{-2}; \& \node{-1}; \\
};
\end{tikzpicture}
\end{itemize}
\end{frame}



\begin{frame}[fragile]{Dictionaries}
A dictionary is an unordered structure of (key, value) pairs.
It is similar to a list but the the objects are accessed by \alert{keys} instead of \alert{index}.

\vfill

One indicates dictionaries by square brackets:

\begin{python}
homework = {'Anna': True, 'Kerstin': False}

homework['Anna'] # True
# changing a value:
homework['Kerstin'] = True
# deleting an item
del homework['Anna']
\end{python}
\vfill
\alert{Note:} Any \emph{immutable} (unchangable) object can be used as a key, e.g. strings, numbers, Booleans, etc. are fine, but not lists.
\vfill

\end{frame}

\begin{frame}[fragile]{Looping through Dictionaries}
A dictionary is an object with the following useful mehods:\\
\vfill
 \hfill \pyth!keys!, \pyth!items!, \pyth!values! \hfill \,
\vfill
By default a dictionary is considered as a list of \alert{keys}:
\begin{python}
for key in homework:  # or homework.keys()
  print(f'{k} {homework[key]}')
\end{python}

\pause
One may also use \pyth!items! to loop through keys and values:
\begin{python}
for key, value in homework.items():
  print('{key} {value}')
\end{python}

\pause
You can also print just the values by using the \pyth!values! method:
\begin{python}
for value in homework.values():
  print('{value}')
\end{python}

\end{frame}

\begin{frame}[fragile]{Dictionaries in this course}
Dictionaries are mainly used for
\vfill
\begin{itemize}
\item providing functions with arguments in a compact way (see Unit 4).
\item to collect options of a method:\\ \pyth!{'tol':1.e-3,'step size': 0.1, 'maxit':1000}!
\end{itemize}
\vfill

\end{frame}

\begin{frame}[fragile]{Sets}
Sets are containers which mimic sets in their mathematical sense:
\begin{definition}
A \alert{set} is a collection of well defined and distinct objects, considered as an object in its own right.
{\footnotesize (Wikipedia)}
\end{definition}
\vfill
\begin{python}
fruitbasket = set(['apple', 'pear', 'banana'])
\end{python}
\vfill
The most important operation is \pyth!in!, meaning $\in$ (is element of):
\begin{python}
'plum' in fruitbasket  # returns False
'pear' in fruitbasket  # returns True
\end{python}
\end{frame}

\begin{frame}[fragile]{Operations on Sets}
Operations on sets are $A \cap B, A \cup B$ and $A / B$ (intersection, union and relative complement):
\begin{python}
fruitbasket = set(['apple', 'pear', 'banana'])
rotten = set(['pear'])
exotic = set(['mango', 'kiwi'])
emptyset = set([])        # the empty set
# Are all fruits edible (non-rotten)?
emptyset == fruitbasket.intersection(rotten)
# Add some diversity
extendedbasket = fruitbasket.union(exotic)
# The edible fruits
goodfruits = fruitbasket - rotten # rel complement
\end{python}
The following syntax can also be used:
\begin{python}
  exotic = {'mango', 'kiwi'}
\end{python}
Compare with dictionaries: \pyth!{'a':'mango', 'b':'kiwi'}!
\end{frame}

\begin{frame}[fragile]{No duplicate elements}
``...distinct objects ...'' (see mathematical definition)\\
\vfill
This is reflected in Python by
\begin{python}
set({'apple', 'apple', 'pear'})
 == set({'apple', 'pear'}) # True
\end{python}
\vfill
\end{frame}




\begin{frame}[fragile]{Tuples}
\begin{defblock}
A \alert{tuple} is an \alert{immutable} list. Immutable means that it cannot be modified.
\end{defblock}

\begin{example}
\begin{python}
my_tuple = (1, 2, 3) # our first tuple!
my_tuple = 1, 2, 3   # parentheses not required
len(my_tuple) # 3, same method as for lists

my_tuple[0] = 'a' # error! tuples are immutable

singleton = 1, # note the comma
singleton = (1,) # preferred for readability
len(singleton) # 1
\end{python}
\end{example}
\end{frame}

\begin{frame}[fragile]{Packing and unpacking}
One may assign several variables at once by \alert{unpacking} a list or tuple:
\begin{python}
a, b = 0, 1 # a gets 0 and b gets 1
a, b = [0, 1] # exactly the same effect
(a, b) = 0, 1 # same
[a,b] = [0,1] # same thing for lists
\end{python}
\pause
\begin{block}{The swapping trick!}
Use packing and unpacking to \alert{swap} the contents of two variables.
\begin{python}
a, b = b, a
\end{python}
\end{block}
\end{frame}


\begin{frame}[fragile]{Returning Multiple Values}
A function may return several values:
\begin{python}
def argmin(L): # return the minimum and index
  ...
  return minimum, minimum_index     # a tuple

min_info = argmin([1, 2, 0])
min_info[0] # 0
min_info[1] # 2
# often unpacking is used directly
m, i = argmin([1, 2, 0]) # m is 0, i is 2
\end{python}
\end{frame}

\begin{frame}[fragile]{A final word on tuples}
\begin{itemize}[<+->]
\item Tuples are \emph{nothing else} than immutable lists
\item In most cases lists may be used instead of tuples
\item The parentheses free notation is nice but \alert{dangerous},
you should \alert{use parentheses when you are not sure}:
\begin{python}
a, b = b, a # the swap trick, equivalent to:
(a, b) = (b, a)
# but
1, 2 == 3, 4 # returns (1, False, 4)
(1, 2) == (3, 4) # returns False
\end{python}
\end{itemize}
\end{frame}




\begin{frame}[fragile]{Conversions}
Containers can be converted to other container types:
\vfill
\begin{tabular}{@{\extracolsep{.5cm}} lll }
  From & To & Command \\
  \hline\\
    List & Tuple & {\pyth!tuple([1, 2, 3])!} \\[1em]
    Tuple & List & {\pyth!list((1, 2, 3))!} \\[1em]
    List, Tuple & Set & {\pyth!set([1,2,3])!, \pyth!set((1,2,3))!} \\[1em]
    Set & List       & {\pyth!list({1,2,3})!} \\[1em]
    Dictionary & List & {\pyth!{'a':4}.values()!} \\[1em]
    List & Dictionary & {--} \\[1em]
  \end{tabular}
\end{frame}

\begin{frame}[fragile]{Summary and Overview}
\vfill
\begin{tabular}{|c||c|c|c|c|} \hline
Type & Access & Order & Duplicate Values & Mutability \\ \hline \hline
List & index & yes & yes & yes \\ \hline
Tuple & index & yes & yes & no \\ \hline
Dictionary & key & no & yes & yes \\ \hline
Set & no & no & no & yes \\ \hline
\end{tabular}
\vfill
\end{frame}

\end{document}
